%! Author = chtompki
%! Date = 1/14/26
% Example LaTeX document for GP111 - note % sign indicates a comment
\documentclass[12pt]{amsart}
% Default margins are too wide all the way around. I reset them here
\setlength{\hoffset}{-.75in}
\setlength{\textwidth}{6.5in}
\setlength{\voffset}{-.5in}
\setlength{\textheight}{9.0in}
\usepackage{amsmath}
\usepackage{amssymb}
\usepackage{amsthm}
\usepackage{enumitem}
\usepackage{setspace}

\newenvironment{solution}
{\begin{proof}[Solution]}
{\end{proof}}

\title{Advanced Linear Algebra\\Homework 1}
\author{Rob Tompkins}
\date{\today}

\begin{document}
    \maketitle
    \date{\today}
    \begin{onehalfspace}

        % Problem P.A.1
        \noindent\textbf{P.A.1} Show that:
        \begin{enumerate}[label=(\alph*)]
            \item $(1 + i)^4 = -4$
            \item $(1 - i)^{-1} - (1 + i)^{-i} = i$
            \item $(i - 1)^{-4} = -1/4$
            \item $10(1 + 3i)^{-1} = 1-3i$
            \item $(\sqrt{3} + i)^3 = 8i$
        \end{enumerate}
        \begin{solution}
            Solution P.A.1
        \end{solution}

        % Problem P.A.3
        \noindent\textbf{P.A.3} Let $z = 3+4i$. Plot and lablel the following pointw in the comples plane: $z$,
        $-z$, $\bar{z}$, $-\bar{z}$, $1/z$, $1/\bar{z}$, $-1/\bar{z}$
        \begin{solution}
            Solution P.A.3
        \end{solution}

        % Problem P.A.6
        \noindent\textbf{P.A.6} Write the following in polar form:
        \begin{enumerate}[label=(\alph*)]
            \item $1 + i$
            \item $(1 + i)^2$
            \item $(1 + i)^4$
        \end{enumerate}
        \begin{solution}
            Solution P.A.6
        \end{solution}

        % Problem P.A.7
        \noindent\textbf{P.A.7} What are the square roots of $1 + i$? Draw a picture of their locations.
        \begin{solution}
            Solution P.A.7
        \end{solution}

        % Problem P.A.11
        \noindent\textbf{P.A.11} If $z \ne 0$, show that $\overline{1/z} = 1/\overline{z}$.
        \begin{solution}
            Solution P.A.11
        \end{solution}

        % Problem P.B.2
        \noindent\textbf{P.B.2} If $p$ is a real polynomial, show that $p(\lambda) = 0$ if
        and only if $p(\bar{\lambda}) = 0$.
        \begin{solution}
            Solution P.B.2
        \end{solution}

        % Problem P.C.3
        \noindent\textbf{P.C.3} If $A\in\mathbf{M}_n(\mathbb{F})$ is invertible, show that
        $A^\top$ is invertible and $(A^{-1})^\top$ is its inverse.
        \begin{solution}
            Solution P.C.3
        \end{solution}

        % Problem P.C.13
        \noindent\textbf{P.C.13} Compute the determinants of:
        $\begin{bmatrix}
             5 & \sqrt{3} + i \\
             \sqrt{3} - i & 4
        \end{bmatrix}$,
        $\begin{bmatrix}
             1 & 2 & 3 \\
             2 & 2 & 3 \\
             3 & 3 & 3
        \end{bmatrix}$, and
        $\begin{bmatrix}
             1 & 1 & 1 \\
             1 & 2 & 2 \\
             1 & 2 & 3
        \end{bmatrix}$
        \begin{solution}
            Solution P.C.13
        \end{solution}

        % Problem P.D.5
        \noindent\textbf{P.D.5} Use mathematical induction to prove that
        \begin{align*}
            1 + z + z^2 + \dots + z^{n-1} = \frac{1 - z^n}{1-z}
        \end{align*}
        \begin{proof}
            Proof P.D.5
        \end{proof}

    \end{onehalfspace}
\end{document}
